\chapter{Reglas, gates y tráfico real}
\label{ch:reglas}

El cliente no sabe nada de DSP por sí mismo. Lo que sabe está en su \textbf{conjunto de reglas}: cada regla
dice qué peticiones cubre, qué parte del mensaje se valida, con qué gates y qué se hace con el resultado.

\section{Una regla}

\begin{lstlisting}[style=sh]
curl -s $ADMINAPI/rules/contract-request
\end{lstlisting}
\begin{lstlisting}[style=json]
{
  "id": "contract-request",
  "priority": 110,
  "listener": "dsp-ingress",
  "match": {
    "method": "POST",
    "path": "/negotiations/{providerPid}/request",
    "body": [ { "pointer": "/@type",
                "one_of": ["ContractRequestMessage", "dspace:ContractRequestMessage"] } ]
  },
  "message": "ContractRequestMessage",
  "verify": { "gates": ["dsp_message", "policy_intake"] },
  "mode": "enforce"
}
\end{lstlisting}

\noindent\begin{tabularx}{\linewidth}{@{}L{4cm}Y@{}}
\toprule
\textbf{Campo} & \textbf{Significado} \\
\midrule
\cod{priority} & Las de valor más bajo se prueban antes. Única por \emph{listener}. \\
\cod{listener} & En qué \emph{listener} aplica. \\
\cod{match.method}, \cod{match.path} & Método y plantilla de ruta. Una variable entre llaves casa con un segmento no vacío. \\
\cod{match.body} & Condiciones opcionales sobre el cuerpo, por puntero JSON. \\
\cod{message} & El tipo de mensaje DSP que la regla espera. \\
\cod{verify.gates} & De uno a cuatro gates. \\
\cod{verify.select} & Puntero JSON al nodo que se valida. Vacío, el cuerpo entero. \\
\cod{mode} & \cod{observe}, por defecto, \cod{shadow} o \cod{enforce}. \\
\cod{validate_response} & Validación pasiva de respuestas de catálogo. \\
\bottomrule
\end{tabularx}

\subsection{Cómo se elige la regla}

Para una petición, los candidatos son las reglas del \emph{listener} cuyo método y ruta casan, en orden de
prioridad. Se elige la primera cuyas condiciones de cuerpo se cumplen. Si no casa ninguna, se aplica
\cod{default_action}.

Se elige \textbf{una sola regla} por mensaje, y cada mensaje se decide con una única revisión del conjunto.

\subsection{Los tres modos}

\noindent\begin{tabularx}{\linewidth}{@{}L{2.6cm}Y Y Y@{}}
\toprule
\textbf{Clasificación} & \textbf{\texttt{observe}} & \textbf{\texttt{shadow}} & \textbf{\texttt{enforce}} \\
\midrule
\ok & reenvía & reenvía & reenvía \\
\warn & reenvía & reenvía y emite evento & reenvía con \cod{VHub-Verdict} y emite evento \\
\viol & reenvía & reenvía y emite evento & \textbf{rechaza} con \cod{400} y emite evento \\
\bottomrule
\end{tabularx}

\begin{ideaClave}
Los tres modos son un camino de despliegue. \cod{observe} no evalúa nada: la regla existe, pero está callada.
\cod{shadow} evalúa y cuenta lo que \emph{se habría} rechazado, sin tocar el tráfico. \cod{enforce} actúa. Una
regla se mueve entre ellos en caliente, y el cambio vale para el mensaje siguiente.
\end{ideaClave}

\subsection{La semilla}

La semilla trae las 16 reglas del enlace DSP 2025-1. Solo \cod{contract-request} viene en \cod{enforce}; las
otras quince, en \cod{observe}.

\begin{longtable}{@{}L{4.5cm}L{6.2cm}L{4.3cm}@{}}
\toprule
\textbf{Regla} & \textbf{Ruta} & \textbf{Gates} \\
\midrule
\endhead
\cod{catalog-request} & \texttt{POST /\allowbreak{}catalog/\allowbreak{}request} & \cod{dsp_message} \\
\cod{dataset-request} & \texttt{GET /\allowbreak{}catalog/\allowbreak{}datasets/\allowbreak{}\{id\}} & \cod{dsp_message} \\
\cod{contract-request-initial} & \texttt{POST /\allowbreak{}negotiations/\allowbreak{}request} & \cod{dsp_message}, \cod{policy_intake} \\
\cod{contract-request} & \texttt{POST /\allowbreak{}negotiations/\allowbreak{}\{providerPid\}/\allowbreak{}request} & \cod{dsp_message}, \cod{policy_intake} \\
\cod{contract-offer-initial} & \texttt{POST /\allowbreak{}negotiations/\allowbreak{}offers} & \cod{dsp_message}, \cod{policy_intake} \\
\cod{contract-offer} & \texttt{POST /\allowbreak{}negotiations/\allowbreak{}\{consumerPid\}/\allowbreak{}offers} & \cod{dsp_message}, \cod{policy_intake} \\
\cod{contract-agreement-in} & \texttt{POST /\allowbreak{}negotiations/\allowbreak{}\{consumerPid\}/\allowbreak{}agreement} & \cod{dsp_message}, \cod{policy_intake} \\
\cod{contract-agreement-out} & La misma, en \cod{dsp-egress} & \cod{dsp_message}, \cod{policy_intake} \\
\cod{agreement-verification} & \texttt{POST …/\allowbreak{}agreement/\allowbreak{}verification} & \cod{dsp_message} \\
\cod{negotiation-event} & \texttt{POST /\allowbreak{}negotiations/\allowbreak{}\{pid\}/\allowbreak{}events} & \cod{dsp_message} \\
\cod{negotiation-termination} & \texttt{POST /\allowbreak{}negotiations/\allowbreak{}\{pid\}/\allowbreak{}termination} & \cod{dsp_message} \\
\cod{transfer-request} & \texttt{POST /\allowbreak{}transfers/\allowbreak{}request} & \cod{dsp_message}, \cod{agreement_admission} \\
\cod{transfer-start} y otras tres & \texttt{POST /\allowbreak{}transfers/\allowbreak{}\{pid\}/\allowbreak{}…} & \cod{dsp_message} \\
\bottomrule
\end{longtable}

\section{El mensaje}

Este es \cod{examples/msgs/m1-forecast.json}: el agregador pide a la distribuidora la oferta del
capítulo~\ref{ch:compilar}, ahora dentro de un mensaje DSP.

\begin{lstlisting}[style=json]
{
  "@context": ["https://w3id.org/dspace/2025/1/context.jsonld",
               "https://flexgrid.example/odrl/profile/metering/context.jsonld"],
  "@type": "ContractRequestMessage",
  "providerPid": "urn:uuid:9c1f7a52-dso-neg-0042",
  "consumerPid": "urn:uuid:4be0d6c1-agg-neg-0017",
  "callbackAddress": "https://aggregator.flexgrid.example/dsp/callback",
  "offer": {
    "@type": "Offer",
    "@id": "urn:uuid:offer-forecast-01",
    "profile": "https://flexgrid.example/odrl/profile/metering",
    "target": "urn:uuid:dataset-qh-load-madrid-2026",
    "permission": [ { "action": "forecastDemand", "constraint": [ … las cinco restricciones … ] } ]
  }
}
\end{lstlisting}

Las dos URL de \cod{@context} son las que el paquete declara con \cod{url}: el contexto de DSP, vendorizado, y
el del perfil, generado. Ninguna se descarga: el cliente las resuelve dentro del paquete.

\section{Preguntar antes de enviar: \texttt{:explain}}

\texttt{POST /\allowbreak{}rules:explain} responde qué haría el proxy con una petición descrita, \textbf{sin efectos}: no
reenvía, no registra y no emite eventos. Recibe \cod{listener}, \cod{method}, \cod{path} y \cod{body}. La
función \cod{ex} lo construye a partir de un fichero de mensaje, y \cod{short} condensa la respuesta.

Con la regla de la semilla:

\begin{lstlisting}[style=sh]
ex m1-forecast | short
\end{lstlisting}
\begin{lstlisting}[style=out]
rule=contract-request state=Active action=reject code=VERDICT_VIOLATION
  gate=dsp_message    class=ok        codes=[] matched=1.0.0
  gate=policy_intake  class=violation codes=['VERDICT_VIOLATION'] matched=None
     - Violation json-schema  missing property 'profile'
     - Info loader  POC_EVAL_UNSUPPORTED_ROLE: shacl artefact profiles/odrl/metering/1.0.0/shapes.ttl is not evaluated
  verdict: violation;gate=policy_intake;code="VERDICT_VIOLATION";rule="contract-request"
\end{lstlisting}

Se rechaza un mensaje que es correcto. El hallazgo dice por qué: falta la propiedad \cod{profile}. Pero el
mensaje la tiene; está dentro de \cod{offer}.

\begin{trampaHabitual}[Una regla, un nodo]
\cod{schema.json} valida \textbf{una política}. La regla de la semilla no tiene \cod{verify.select}, así que el
gate \cod{policy_intake} recibe el mensaje entero, y un \cod{ContractRequestMessage} no es una política: la
política está en \cod{/offer}.

La solución es \texttt{"select": "/\allowbreak{}offer"}. Pero \cod{select} es \textbf{uno por regla} y vale para todos sus
gates, y solo se elige una regla por mensaje. Con \texttt{select: "/\allowbreak{}offer"}, el gate \cod{dsp_message} recibiría
también la oferta en lugar de la envoltura. Hoy no se pueden validar en un mismo mensaje la envoltura con un
esquema y la política con otro: hay que elegir.
\end{trampaHabitual}

Para una petición de contrato, lo que importa es la política. La regla adaptada está en
\cod{examples/rules/contract-request.json}:

\begin{lstlisting}[style=json]
  "verify": { "select": "/offer", "gates": ["policy_intake"] },
  "mode": "enforce"
\end{lstlisting}

\section{Cambiar una regla}

Todo cambio de reglas exige \cod{If-Match} con la revisión actual, que es el \cod{ETag} de \texttt{GET /\allowbreak{}rules}.

\begin{lstlisting}[style=sh]
curl -s -i -X PUT $ADMINAPI/rules/contract-request --data-binary @../rules/contract-request.json | head -1
curl -s -i -X PUT $ADMINAPI/rules/contract-request -H 'If-Match: "1"' \
  --data-binary @../rules/contract-request.json | grep -i -E '^HTTP|^etag'
\end{lstlisting}
\begin{lstlisting}[style=out]
HTTP/1.1 428 Precondition Required

HTTP/1.1 200 OK
Etag: "2"
\end{lstlisting}

Sin la cabecera, \cod{428}. Con una revisión desfasada sería \cod{412}. El conjunto entero se revalida en cada
cambio, y un cambio rechazado deja intacto el conjunto anterior.

\noindent\begin{tabularx}{\linewidth}{@{}C{1.8cm}Y L{5.7cm}@{}}
\toprule
\textbf{Invariante} & \textbf{Comprobación} & \textbf{Código} \\
\midrule
RULE-1 & El conjunto valida contra su esquema. & \cod{RULES_SCHEMA} \\
RULE-2 & Los identificadores son únicos. & \cod{RULES_DUPLICATE_ID} \\
RULE-3 & Las prioridades son únicas por \emph{listener}. & \cod{RULES_DUPLICATE_PRIORITY} \\
RULE-4 & \cod{agreement_admission} solo en peticiones de transferencia; \cod{validate_response} solo en catálogo. & \cod{RULES_GATE_MESSAGE_MISMATCH} \\
RULE-5 & Las reglas de \cod{mgmt-ingress} necesitan \cod{upstreams.management}. & \cod{RULES_UPSTREAM_MISSING} \\
\bottomrule
\end{tabularx}

\medskip
De paso se adaptan las otras dos reglas que usará el capítulo siguiente. \cod{setrule} lee la revisión actual y
hace el \cod{PUT}:

\begin{lstlisting}[style=sh]
setrule contract-agreement-in
setrule transfer-request
\end{lstlisting}
\begin{lstlisting}[style=out]
contract-agreement-in -> 200
transfer-request -> 200
\end{lstlisting}

\noindent\begin{tabularx}{\linewidth}{@{}L{4.7cm}Y@{}}
\toprule
\textbf{Regla} & \textbf{Cambio respecto de la semilla} \\
\midrule
\cod{contract-request} & \texttt{select: /\allowbreak{}offer}, solo \cod{policy_intake}. Ya estaba en \cod{enforce}. \\
\cod{contract-agreement-in} & \texttt{select: /\allowbreak{}agreement}, solo \cod{policy_intake}, pasa a \cod{enforce}. \\
\cod{transfer-request} & Pasa a \cod{enforce}. Conserva \cod{dsp_message} y \cod{agreement_admission} sobre el cuerpo entero. \\
\bottomrule
\end{tabularx}

\section{Seis mensajes}

Con la regla adaptada, estos son los veredictos a 3 de septiembre, con la 1.0.0 vigente:

\begin{lstlisting}[style=sh]
for m in m1-forecast m2-profiling m3-research m4-invalid m5-foreign-context m6-no-callback; do
  echo "--- $m"; ex $m | short
done
\end{lstlisting}
\begin{lstlisting}[style=out]
--- m1-forecast
rule=contract-request state=Active action=forward code=None
  gate=policy_intake  class=ok        codes=[] matched=1.0.0
--- m2-profiling
rule=contract-request state=Active action=forward code=None
  gate=policy_intake  class=ok        codes=[] matched=1.0.0
--- m3-research
rule=contract-request state=Active action=reject code=VERDICT_VIOLATION
  gate=policy_intake  class=violation codes=['VERDICT_VIOLATION'] matched=None
     - Violation json-schema /permission/0/constraint/0 validation failed
--- m4-invalid
rule=contract-request state=Active action=reject code=VERDICT_VIOLATION
  gate=policy_intake  class=violation codes=['VERDICT_VIOLATION'] matched=None
     - Violation json-schema /permission/0/action validation failed
     - Violation json-schema /permission/1/constraint/0 validation failed
--- m5-foreign-context
rule=contract-request state=Active action=reject code=CONTEXT_NOT_IN_PACKAGE
  gate=policy_intake  class=violation codes=['CONTEXT_NOT_IN_PACKAGE'] matched=None
     - Violation loader  @context https://schema.org/ is not in the package
--- m6-no-callback
rule=contract-request state=Active action=forward code=None
  gate=policy_intake  class=ok        codes=[] matched=1.0.0
\end{lstlisting}

En todas las respuestas aparece además la línea informativa \cod{POC_EVAL_UNSUPPORTED_ROLE}, omitida aquí.

\noindent\begin{tabularx}{\linewidth}{@{}L{3.9cm}Y C{1.9cm}@{}}
\toprule
\textbf{Mensaje} & \textbf{Qué lleva} & \textbf{Resultado} \\
\midrule
\cod{m1-forecast} & Previsión de demanda con cinco restricciones válidas. & \ok \\
\cod{m2-profiling} & Perfilado de consumidores. En la 1.0.0 es una acción activa. & \ok \\
\cod{m3-research} & Finalidad \cod{enr:research}, que la 1.0.0 no conoce. & \viol \\
\cod{m4-invalid} & Acción \cod{odrl:sell} y \texttt{minAggregation lt 5}. Dos fallos, dos hallazgos. & \viol \\
\cod{m5-foreign-context} & Un tercer contexto, \cod{https://schema.org/}. & \viol \\
\cod{m6-no-callback} & Sin \cod{callbackAddress}, obligatorio en DSP. & \ok \\
\bottomrule
\end{tabularx}

\medskip
Tres comentarios.

\textbf{m3} es una violación \emph{hoy}. El capítulo siguiente publica la 1.1.0 y el mismo mensaje pasa.

\textbf{m5} no se rechaza por su política, que es válida, sino por citar un contexto que el paquete no contiene.
El cliente nunca descarga un contexto para validar: si lo hiciera, quien controlara esa URL podría cambiar el
significado de un mensaje después de firmado el vocabulario. Esta comprobación la dispara el artefacto
\cod{shacl} del \emph{binding}: con solo \cod{json-schema}, los contextos no se miran.

\textbf{m6} pasa aunque le falta un campo obligatorio de DSP. Es el precio de la trampa anterior: esta regla ya
no evalúa \cod{dsp_message}.

\section{Tráfico de verdad}

\cod{row} envía un mensaje por \cod{dsp-ingress} e imprime el estado y la cabecera de veredicto.

\begin{lstlisting}[style=sh]
curl -s -i -X POST "$INGRESS/negotiations/$NEG/request" -H 'Content-Type: application/json' \
  --data-binary @../msgs/m1-forecast.json
\end{lstlisting}
\begin{lstlisting}[style=out]
HTTP/1.1 200 OK
Content-Length: 90
Content-Type: application/json

{"@type":"Ack","received":"POST /dsp/negotiations/urn:uuid:9c1f7a52-dso-neg-0042/request"}
\end{lstlisting}

Es la respuesta del conector, sin tocar. El conector de laboratorio recuerda lo que recibe:

\begin{lstlisting}[style=sh]
curl -s $CONNECTOR/_received
\end{lstlisting}

El cuerpo que aparece ahí tiene los mismos 1627 bytes que el fichero enviado.

Ahora el mensaje inválido:

\begin{lstlisting}[style=sh]
curl -s -i -X POST "$INGRESS/negotiations/$NEG/request" -H 'Content-Type: application/json' \
  --data-binary @../msgs/m4-invalid.json
\end{lstlisting}
\begin{lstlisting}[style=out]
HTTP/1.1 400 Bad Request
Content-Type: application/json
VHub-Verdict: violation;gate=policy_intake;code="VERDICT_VIOLATION";rule="contract-request"

{"@context":["https://w3id.org/dspace/2025/1/context.jsonld"],
 "@type":"ContractNegotiationError","code":"VERDICT_VIOLATION",
 "consumerPid":"urn:uuid:4be0d6c1-agg-neg-0017","providerPid":"urn:uuid:9c1f7a52-dso-neg-0042",
 "reason":["validation failed","validation failed"]}
\end{lstlisting}

El conector no ha recibido nada. La respuesta es el mensaje de error que DSP define para ese intercambio, un
\cod{ContractNegotiationError} con los identificadores del proceso copiados de la petición. Un conector al otro
lado lo entiende sin saber que existe Alexandria.

\cod{reason} lleva un elemento por hallazgo, dos en este caso. Los hallazgos informativos, como el aviso de que
el artefacto \cod{shacl} no se ejecuta, no se envían a la contraparte: quedan en \cod{:explain} y en el evento.

\subsection{La cabecera \texttt{VHub-Verdict}}

Es una lista de campos estructurados (RFC~9651), con un elemento por cada gate en \cod{enforce} cuya
clasificación no es \ok:

\begin{lstlisting}[style=out]
VHub-Verdict: violation;gate=policy_intake;code="VERDICT_VIOLATION";rule="contract-request"
VHub-Verdict: warning;gate=policy_intake;code="VERDICT_NOT_IN_HEAD";rule="contract-request";version="1.1.0";sunset=@1793613600
\end{lstlisting}

\section{Eventos y métricas}

Cada decisión que merece contarse se añade al fichero de eventos:

\begin{lstlisting}[style=sh]
tail -n 1 var/client/events.jsonl
\end{lstlisting}
\begin{lstlisting}[style=json]
{"type":"vhub.verdict","ts":"2026-09-03T12:00:00Z","listener":"dsp-ingress",
 "rule":"contract-request","rules_revision":2,"message":"ContractRequestMessage",
 "gate":"policy_intake","mode":"enforce","state":"Active",
 "classification":"violation","action":"reject","evaluated":["1.0.0"],"matched":null,
 "codes":["VERDICT_VIOLATION"],
 "findings":[
   {"slot":"odrl","source":"json-schema","severity":"Violation","path":"/permission/0/action","message":"validation failed"},
   {"slot":"odrl","source":"json-schema","severity":"Violation","path":"/permission/1/constraint/0","message":"validation failed"},
   … ]}
\end{lstlisting}

El evento dice qué regla, en qué revisión, contra qué versiones se evaluó, cuál aceptó el mensaje y, por
hallazgo, el slot, el artefacto y el puntero al nodo culpable.

\noindent\begin{tabularx}{\linewidth}{@{}L{3cm}Y@{}}
\toprule
\textbf{Tipo} & \textbf{Cuándo} \\
\midrule
\cod{vhub.verdict} & Un gate produce una clasificación que hay que contar. \\
\cod{vhub.sync} & Falla una sincronización. También lleva los avisos del evaluador. \\
\cod{vhub.state} & Cambia el estado, o el desfase de reloj supera el tolerado. \\
\cod{vhub.agreement} & Se programa la revocación de la versión de un acuerdo capturado. \\
\cod{vhub.rules} & Cambia el conjunto de reglas. \\
\cod{vhub.audit} & Se produce un informe de auditoría pasiva. \\
\cod{vhub.gc} & Pasa el recolector de paquetes. \\
\bottomrule
\end{tabularx}

\medskip
Y las métricas, en formato OpenMetrics:

\begin{lstlisting}[style=sh]
curl -s $ADMINAPI/metrics | grep -v '^#'
\end{lstlisting}
\begin{lstlisting}[style=out]
vhub_classifications_total{rule="contract-request",gate="policy_intake",version="",result="violation"} 1
vhub_classifications_total{rule="contract-request",gate="policy_intake",version="1.0.0",result="ok"} 1
vhub_listing_age_seconds 0
vhub_request_latency_seconds_sum 0.058964584
vhub_request_latency_seconds_count 2
\end{lstlisting}

\cod{vhub_listing_age_seconds} es la métrica sobre la que poner una alerta: si crece, el cliente ha dejado de
hablar con el hub, y cuando pase de \cod{ttl} entrará en \cod{Degraded}.

\begin{retos}
\begin{enumerate}[leftmargin=*, itemsep=2pt]
  \item Pasa \cod{contract-request} a \cod{shadow} y envía \cod{m4-invalid}. ¿Qué recibe el emisor? ¿Qué queda
        en el fichero de eventos?
  \item Añade con \texttt{POST /\allowbreak{}rules} una copia de una regla existente. ¿Qué invariante salta?
  \item Cambia \cod{default_action} a \cod{reject} en la configuración, reinicia el cliente y envía una petición
        a una ruta sin regla.
\end{enumerate}
\end{retos}
